% Generated by task tex-libraries from dossiers/lib-others.md, section D.
\begingroup\sloppy
\section{Mathlib and Lean itself: what the statements unfold to, and the trusted base}\label{sec:gen-lib-lean}

\subsection{Mathlib notions in the proof system's statements}\label{sec:gen-lib-lean-mathlib}

\Cref{tab:gen-lib-mathlib} lists the Mathlib (and core) notions counted in the statements, not
the proofs, of \file{LeanerVM/Protocol/} at \code{b435631} and in the library objects of the
previous sections. Definitions are quoted from Mathlib at the old pin \code{v4.33.1}
(\code{0df444a3}, read with \code{git show 0df444a3:<path>}) and core at Lean \code{v4.33.1}.

{\small
\begin{longtable}{@{}p{0.15\linewidth}p{0.17\linewidth}p{0.34\linewidth}p{0.26\linewidth}@{}}
\caption{Mathlib and core notions the proof system's statements unfold to.}\label{tab:gen-lib-mathlib}\\
\toprule
Notion & What it is & Definition (old pin) & Where the proof system meets it \\
\midrule
\endfirsthead
\toprule
Notion & What it is & Definition (old pin) & Where the proof system meets it \\
\midrule
\endhead
\bottomrule
\endfoot
\lean{ℝ≥0}, \lean{ℝ≥0∞} & non-negative reals, and with $\infty$
  & \lean{def NNReal := { r : ℝ // 0 ≤ r }} (\file{Mathlib/Data/NNReal/}\allowbreak\file{Defs.lean:58});
    \lean{def ENNReal := WithTop ℝ≥0} (\file{Mathlib/Data/ENNReal/}\allowbreak\file{Basic.lean:101})
  & error bounds: \lean{piopError : … → ℝ≥0} (\file{Spine/Compose.lean:149}); every
    \lean{Pr[…]} is an \lean{ℝ≥0∞} \\
\lean{PMF}, \lean{uniformOf}\allowbreak\lean{Fintype} & probability mass function; the uniform one
  & \lean{def PMF (α) := { f : α → ℝ≥0∞ // HasSum f 1 }}
    (\file{Probability/}\allowbreak\file{ProbabilityMassFunction/}\allowbreak\file{Basic.lean:46-47});
    \lean{uniformOfFintype α := uniformOfFinset Finset.univ …}
    (\file{Probability/Distributions/}\allowbreak\file{Uniform.lean:285-286})
  & the semantics of every coin (\cref{sec:gen-lib-vcvio-dist}) \\
\lean{Set α} & a predicate on \lean{α}
  & core/Mathlib \lean{Set α := α → Prop}
  & ArkLib relations: \lean{M3Rel I : Set ((I.Stmt × …) × Column I.μ)}
    (\file{Spine/Instance.lean:230}), \lean{Seam.of} (\file{Seams.lean:166}), language sets in
    \file{Spine/Phase.lean} \\
\lean{Fintype α}, \lean{Fintype.card} & a finite type with its list of elements; its size
  & \lean{class Fintype (α) where elems : Finset α; complete : ∀ x, x ∈ elems}
    (\file{Data/Fintype/}\allowbreak\file{Defs.lean:57-61}); \lean{def card (α) [Fintype α] : ℕ := (@univ α _).card}
    (\file{Data/Fintype/}\allowbreak\file{Card.lean:39-41})
  & the size of the challenge space, \lean{card_E} (\cref{sec:gen-l0-catalogue}); the counting
    bounds (\cref{sec:gen-l0-bounds}); \file{PublicInput.lean} \\
\lean{Finset.sum} ($\sum$), \lean{Finset.prod} & finite sums and products
  & Mathlib \lean{Finset.sum s f} (fold over the underlying multiset)
  & stacking and claim weights: \file{Seams.lean:113}
    (\lean{∑ i : Fin (2 ^ μ), W.onCube.get i * ofK (q.values.get i)}), \file{Stack},
    \file{Padding}, \file{ClaimWeights}, \file{FixedColumns} \\
\lean{List.Perm} (\lean{~}) & one list is a reordering of the other (multiset equality)
  & core \lean{inductive Perm : List α → List α → Prop} with \lean{nil}, \lean{cons},
    \lean{swap}, \lean{trans} (\file{Init/Data/List/}\allowbreak\file{Basic.lean:1858-1874}, Lean \code{v4.33.1})
  & bus balance: \lean{def Balanced (q) : Prop := (I.tuples q .push).Perm (I.tuples q .pull)}
    (\file{Spine/Instance.lean:201}), leanISA's \lean{BalancedPair}
    (\cref{sec:gen-lib-clean-workaround}) \\
\lean{MvPolynomial σ R} & Mathlib's multivariate polynomials (non-computable, finitely
  supported coefficient maps)
  & \lean{abbrev MvPolynomial (σ R) := AddMonoidAlgebra R (σ →₀ ℕ)}
    (\file{Algebra/MvPolynomial/}\allowbreak\file{Basic.lean:82-83})
  & reached through \lean{fromCMvPolynomial} (\cref{sec:gen-lib-comppoly-mv}) for
    \lean{totalDegree} and Schwartz--Zippel \\
\lean{Polynomial}, \lean{AdjoinRoot}, \lean{Irreducible} & univariate polynomials, the quotient
  $R[X]/(f)$, irreducibility
  & \lean{def AdjoinRoot (f : R[X]) := Polynomial R ⧸ (span {f} : Ideal R[X])}
    (\file{RingTheory/}\allowbreak\file{AdjoinRoot.lean:62})
  & the \emph{definitions} of $\K$ and $\E$ (\cref{sec:gen-lib-comppoly-fields}), through the
    bijections \lean{toQuot} \\
\end{longtable}}

A reader who knows these notions from mathematics can take them at face value: each is a short
definition whose meaning is the textbook one, and Mathlib's theory about them is proved.

\subsection{The trusted base of any Lean development}\label{sec:gen-lib-lean-tcb}

\begin{define}{The kernel and the three standard axioms}
\textbf{The kernel} is Lean's type checker for its dependent type theory. Every declaration,
however produced (by hand, by tactics, by automation), is re-checked by it; tactics are
untrusted. \textbf{Three standard axioms} are assumed: \lean{propext} (logically equivalent
propositions are equal), \lean{Classical.choice} (choice; gives excluded middle) and
\lean{Quot.sound} (elements related by a relation have equal classes in its quotient; gives function
extensionality). They are the axioms of Mathlib and are consistent relative to standard set
theory.
\end{define}
(The probe \code{NumeralHazard} shows the kernel rejecting two ill-typed proof terms that
\lean{simp} produced: \cref{f:lib-simp-misreads-k}.)

\paragraph{What would enlarge it.} \lean{sorryAx} (an admitted proof: it proves anything);
\lean{native_decide}, which adds the axiom \lean{Lean.ofReduceBool} and so trusts the compiler,
the interpreter and every \lean{@[implemented_by]}/\lean{@[extern]} replacement of a definition
by native code; any user \lean{axiom}.

\paragraph{How closed claims are checked, and what each way trusts.}
\begin{description}
\item[\lean{decide}] The elaborator evaluates the \lean{Decidable} instance, and the kernel
  re-checks the resulting term \lean{of_decide_eq_true (Eq.refl true)} by its own reduction.
  Trusted: the kernel.
\item[\lean{decide +kernel}] Skips the elaborator's evaluation and asks the kernel directly (the
  form leanerVM uses for $\K$ powers, \file{Parameters/Generator.lean:281-299}). Trusted: the
  kernel.
\item[\texttt{\#guard e}] Evaluates \lean{e} by the \textbf{compiler or interpreter} and fails the
  build if it is not \lean{true}. It proves nothing to the kernel and trusts compiled code,
  including \lean{@[csimp]} replacements (CompPoly's \lean{mul_eq_mulTbl} swaps \lean{Ext}
  multiplication for a table version; that equation is itself proved,
  \file{Extension/Defs.lean:267}) and core's native \lean{Nat}/\lean{BitVec} code. leanerVM uses
  \texttt{\#guard} only in tests, as executable evidence.
\item[\lean{native_decide}] Decides the claim by compiled code and records the axiom
  \lean{Lean.ofReduceBool}: it trusts the compiler, the interpreter and every native replacement
  of a definition. Banned in leanerVM (\cref{sec:gen-lib-lean-validation}).
\item[\lean{rfl}] \lean{example … := rfl} is kernel definitional equality. Trusted: the kernel.
\end{description}

\subsection{What leanerVM's validation enforces (at \texorpdfstring{\code{b435631}}{b435631})}\label{sec:gen-lib-lean-validation}

\begin{itemize}
\item \file{scripts/audit-lean.sh:8-21}: a lexical ban on \lean{axiom}, \lean{sorry},
  \lean{admit}, \lean{unsafe}, \lean{native_decide} in \file{LeanerVM}, \file{LeanerVM.lean} and
  \file{tests}, plus \file{check-lean-options.py} (no local override of linter or
  implicit-variable options).
\item \file{scripts/audit-axioms.lean:17-29}: for \textbf{every} constant whose name starts with
  \lean{LeanerVM} or \lean{LeanerVMTests}, \lean{collectAxioms} (recursive, through all imported
  libraries) must return only \lean{propext}, \lean{Classical.choice}, \lean{Quot.sound}.
\end{itemize}
\begin{leancode}
for (name, _) in env.constants.toList do
  if (`LeanerVM).isPrefixOf name || (`LeanerVMTests).isPrefixOf name then
    count := count + 1
    for ax in ← collectAxioms name do
      unless #[`propext, `Classical.choice, `Quot.sound].contains ax do
        unexpected := unexpected.push (name, ax)
\end{leancode}
\src{\file{scripts/audit-axioms.lean:17-29}, leanerVM \code{b435631}}
\begin{itemize}
\item \file{scripts/test-axiom-audit.py:28-45} re-runs the audit with a planted
  \lean{axiom AuditNegative.hidden : False} used by one production and one test theorem, and
  requires the audit to reject both (a negative control). CI runs the same audit through
  \code{lean-action} (\code{axiom-audit-root: LeanerVM}, \file{.github/workflows/ci.yml:40-43}) and
  the script (\file{:44-45}); \file{docs/ci.md:8-9} marks both workflows required.
\item \file{scripts/validate.sh}: the above plus the import DAG, layer and documentation checks,
  \code{lake build}, \code{lake test}, and the test root with warnings as errors.
\end{itemize}

\paragraph{Consequence.} A leanerVM theorem cannot silently rest on one of ArkLib's admitted
theorems (ArkLib's \lean{sorryAx} would appear in \lean{collectAxioms}), nor on
\lean{native_decide} anywhere in its dependency cone. What remains trusted is the kernel, the
three axioms, and the \emph{meaning of the definitions} in the statements
(\cref{sec:gen-lib-vcvio-trust,sec:gen-lib-comppoly,sec:gen-lib-clean,sec:gen-lib-lean-mathlib}).
\texttt{\#guard} evidence in tests is compiled-code evidence, never proof. (The audit covers the two
project namespaces; an instance declared outside them would escape it. None of Layer 0's is:
\lean{LeanerVM.Protocol.instSampleableTypeK} and the others are in the namespace, per the probe
\code{Layer0}'s output.)

At \code{144c5aa} the pins are Mathlib and Lean \code{v4.34.1}; \file{docs/dependencies.md} at
that revision describes the upgrade.

\par\endgroup